% Luc Destombe --- licence CC BY-NC-SA 4.0
% https://creativecommons.org/licenses/by-nc-sa/4.0/deed.fr
% Fichier autonome : compiler avec pdflatex, deux fois (signets, nombre de pages).
\documentclass[11pt,a4paper]{article}
\makeatletter
\newif\ifeleve
\elevefalse
\elevetrue

\newif\iflycee@palatino
\lycee@palatinotrue

\RequirePackage[utf8]{inputenc}
\RequirePackage[T1]{fontenc}
\RequirePackage[french]{babel}
\RequirePackage{etoolbox}

\newcommand{\ShorthandsOff}{\@ifundefined{shorthandoff}{}{\shorthandoff{;:!?}}}
\newcommand{\ShorthandsOn}{\@ifundefined{shorthandon}{}{\shorthandon{;:!?}}}
\ShorthandsOff
\AtBeginDocument{\ShorthandsOn}
\AtBeginEnvironment{tikzpicture}{\ShorthandsOff}

\RequirePackage{geometry}
\RequirePackage{fancyhdr}
\RequirePackage{lastpage}
\RequirePackage{multicol}
\RequirePackage{array}
\RequirePackage{enumitem}
\RequirePackage{xcolor}
\RequirePackage{needspace}

\RequirePackage{amsmath,amssymb}
\RequirePackage{mathpazo}
\DeclareMathSizes{10.95}{10.95}{8}{5.5}
\begingroup\catcode`\_=\active \catcode`\^=\active
\gdef_#1{\sb{\scriptscriptstyle #1}}
\gdef^#1{\sp{\scriptscriptstyle\lycee@moinsepais #1}}
\endgroup
\newcommand\lycee@moins{\mathbin{%
  \setbox\z@\hbox{$\m@th\scriptscriptstyle\lycee@moinsorig$}%
  \dimen@=.55\wd\z@ \advance\dimen@-.5pt
  \hbox to.75\wd\z@{\hss$\m@th\scriptscriptstyle\vcenter{\hbox{%
    \tikz\draw[line width=.5pt,line cap=round](0,0)--(\the\dimen@,0);}}$\hss}}}
\begingroup\lccode`\~=`\- \lowercase{\endgroup\let~}\lycee@moins
\newcommand\lycee@moinsepais{\mathcode`\-="8000 }
\AtBeginDocument{\mathchardef\lycee@moinsorig=\mathcode`\-\relax}
\AtBeginDocument{\catcode`\_=12 \mathcode`\_="8000
  \catcode`\^=12 \mathcode`\^="8000 }
\newcommand\lycee@exposants{%
  \fontdimen13\textfont2=.48\fontdimen6\textfont2
  \fontdimen14\textfont2=.48\fontdimen6\textfont2
  \fontdimen15\textfont2=.48\fontdimen6\textfont2
  \fontdimen13\scriptfont2=.48\fontdimen6\scriptfont2
  \fontdimen14\scriptfont2=.48\fontdimen6\scriptfont2
  \fontdimen15\scriptfont2=.48\fontdimen6\scriptfont2}
\every@math@size\expandafter{\the\every@math@size\lycee@exposants}
\newlength{\retraitcalculs}
\setlength{\retraitcalculs}{2em}

\RequirePackage{tikz}
\usetikzlibrary{arrows.meta,calc,babel,shapes.misc}
\RequirePackage[most]{tcolorbox}
\tcbuselibrary{skins,breakable}

\definecolor{coulDef}{HTML}{1F4E79}
\definecolor{coulProp}{HTML}{7030A0}
\definecolor{coulRem}{HTML}{C55A11}
\definecolor{coulEx}{HTML}{2E7D32}
\definecolor{coulExo}{HTML}{B03A2E}
\definecolor{coulPart}{HTML}{1F4E79}
\definecolor{coulMeth}{HTML}{1B6E4F}
\definecolor{coulCourbe}{HTML}{0B5ED7}
\definecolor{coulCourbeB}{HTML}{C00000}
\definecolor{coulCourbeC}{HTML}{2E7D32}
\definecolor{coulGrille}{HTML}{8C8C8C}
\definecolor{coulRep}{HTML}{1B6E4F}
\definecolor{coulBar}{HTML}{2E7D32}

\newcommand{\R}{\mathbb{R}}
\newcommand{\Rs}{\mathbb{R}^{*}}
\renewcommand{\le}{\leqslant}
\renewcommand{\ge}{\geqslant}
\newcommand{\vect}[1]{\overrightarrow{#1}}
\newcommand{\norme}[1]{\left\lVert #1 \right\rVert}
\newcommand{\intervalle}[2]{\left[#1\,;#2\right]}
\RequirePackage{cancel}
\renewcommand{\CancelColor}{\color{coulExo}}
\newcommand{\fois}[1]{\times{\color{coulExo}#1}}
\newcommand{\barre}[1]{\cancel{#1}}

\tikzset{grille/.style={coulGrille,line width=0.4pt}}
\newcommand{\quadrillage}[4]{\draw[grille,step=1] (#1,#2) grid (#3,#4);}
\tikzset{
  croix/.style={cross out,draw,line width=0.6pt,inner sep=0pt,minimum size=4pt},
  croixdroite/.style={croix,rotate=45,line width=0.8pt},
}
\newcommand{\pt}[3]{\path (#1) node[croix]{} node[#2,font=\small,inner sep=2.5pt]{$#3$};}

\newcommand{\lycee@repere}[4]{%
  \draw[grille,step=1] (#1,#2) grid (#3,#4);
  \draw[-{Stealth[length=2mm]},thick] (#1,0)--({#3+0.4},0) node[below left=-1pt]{$x$};
  \draw[-{Stealth[length=2mm]},thick] (0,#2)--(0,{#4+0.4}) node[below left=-1pt]{$y$};
  \node[below left,font=\scriptsize,inner sep=1.5pt] at (0,0) {$0$};}
\newcommand{\lycee@gradx}[1]{\foreach \k/\l in {#1}{%
  \draw (\k,0.07)--(\k,-0.07) node[below,font=\scriptsize,inner sep=1.5pt]{$\l$};}}
\newcommand{\lycee@grady}[1]{\foreach \k/\l in {#1}{%
  \draw (0.07,\k)--(-0.07,\k) node[left,font=\scriptsize,inner sep=1.5pt]{$\l$};}}
\newcommand{\lycee@intff}[2]{\left[#1\,;#2\right]}
\newcommand{\lycee@intfo}[2]{\left[#1\,;#2\right[}
\newcommand{\lycee@intof}[2]{\left]#1\,;#2\right]}
\newcommand{\lycee@intoo}[2]{\left]#1\,;#2\right[}
\newcommand{\lycee@ens}[1]{\left\{#1\right\}}
\AtBeginDocument{%
  \providecommand{\repere}{\lycee@repere}%
  \providecommand{\gradx}{\lycee@gradx}%
  \providecommand{\grady}{\lycee@grady}%
  \providecommand{\Cf}{\mathcal{C}\sb{\scriptscriptstyle f}}%
  \providecommand{\Cg}{\mathcal{C}\sb{\scriptscriptstyle g}}%
  \providecommand{\intff}{\lycee@intff}%
  \providecommand{\intfo}{\lycee@intfo}%
  \providecommand{\intof}{\lycee@intof}%
  \providecommand{\intoo}{\lycee@intoo}%
  \providecommand{\ens}{\lycee@ens}}

\newlist{questions}{enumerate}{3}
\setlist[questions,1]{label=\arabic*\textdegree),leftmargin=*,itemsep=2pt,topsep=3pt}
\setlist[questions,2]{label=\alph*),leftmargin=*,itemsep=1pt,topsep=2pt}
\setlist[questions,3]{label=\roman*),leftmargin=*,itemsep=1pt,topsep=1pt}
\setlist[itemize]{label=\textbullet,leftmargin=*,itemsep=2pt,topsep=3pt}

\newcommand{\besoinplace}[1]{\par\penalty\@M\begingroup
  \dimen@=#1\relax\dimen@ii=\pagegoal\advance\dimen@ii-\pagetotal
  \ifdim\dimen@>\dimen@ii\ifdim\dimen@ii>\z@\vfil\fi\break\fi\endgroup}

\newcommand{\lycee@pied}{\thepage/\pageref{LastPage}}
\newcommand{\entete}[2]{%
  \pagestyle{fancy}\fancyhf{}%
  \fancyhead[L]{\footnotesize\color{coulPart}\textbf{#1}}%
  \fancyhead[R]{\footnotesize\color{coulPart}\textbf{#2}}%
  \fancyfoot[C]{\footnotesize\lycee@pied}%
  \renewcommand{\headrulewidth}{0.4pt}%
  \renewcommand{\headrule}{\hbox to\headwidth{%
    \color{coulPart!40}\leaders\hrule height \headrulewidth\hfill}}}

\RequirePackage{ccicons}
\newcommand{\lycee@licence}{{\scriptsize\color{gray}%
  \copyright\ Luc Destombe --- \ccbyncsa\ CC BY-NC-SA 4.0}}
\newif\iflycee@licence \lycee@licencetrue
\newcommand{\sanslicence}{\lycee@licencefalse}
\AtBeginDocument{\iflycee@licence\AddToHookNext{shipout/foreground}{%
  \put(0,0){\rlap{\kern\dimexpr1in+\hoffset+\oddsidemargin+\textwidth\relax
    \llap{\raisebox{-\dimexpr1in+\voffset+\topmargin+\headheight+\headsep
      +\textheight+\footskip\relax}{\lycee@licence}}}}}\fi}

\newcommand{\formulesserrees}{%
  \setlength{\abovedisplayskip}{3pt}\setlength{\belowdisplayskip}{3pt}%
  \setlength{\abovedisplayshortskip}{2pt}\setlength{\belowdisplayshortskip}{3pt}}

\setlength{\parindent}{0pt}

\RequirePackage[implicit=false,hidelinks,pdfencoding=auto,psdextra,bookmarksopen,
  bookmarksopenlevel=1,pdfcreator={},pdfproducer={}]{hyperref}
\RequirePackage{bookmark}
\hypersetup{pdfauthor={Luc Destombe},pdfkeywords={CC BY-NC-SA 4.0}}
\newcounter{lycee@signet}
\newcommand{\signet}[3][]{\stepcounter{lycee@signet}%
  \edef\lycee@dest{\if\relax\detokenize{#1}\relax
    signet.\arabic{lycee@signet}\else#1\fi}%
  \hypertarget{\lycee@dest}{}\bookmark[level=#2,dest=\lycee@dest]{#3}}
\pdfstringdefDisableCommands{%
  \def\R{R}\def\vect#1{#1}\def\Cf{Cf}\def\Cg{Cg}\def\fois#1{×#1}%
  \def\barre#1{#1}\def\intervalle#1#2{[#1;#2]}\def\intff#1#2{[#1;#2]}%
  \def\textsuperscript#1{#1}\def\dfrac#1#2{#1/#2}\def\frac#1#2{#1/#2}%
  \def\vec#1{#1⃗}}%
\geometry{top=2cm,bottom=1cm,left=1cm,right=1cm,headheight=14pt,footskip=0.6cm}

\RequirePackage{pgfplots}
\pgfplotsset{compat=1.18}
\RequirePackage{tkz-tab}

\newcounter{partiecours}
\renewcommand{\thepartiecours}{\Roman{partiecours}}
\newcommand{\partiecours}[1]{%
  \refstepcounter{partiecours}%
  \Needspace*{8\baselineskip}%
  \bigskip
  \noindent\begin{tcolorbox}[enhanced,colback=coulPart,colframe=coulPart,
      boxrule=0pt,arc=2pt,left=8pt,right=8pt,top=4pt,bottom=4pt,
      before skip=6pt,after skip=10pt]
    \color{white}\bfseries\large \thepartiecours{} --- #1%
    \signet[partie-\arabic{partiecours}]{1}{\thepartiecours{} --- #1}
  \end{tcolorbox}%
  \addcontentsline{toc}{section}{\thepartiecours{} --- #1}%
}

\newtcolorbox{boitecours}[3][]{%
  enhanced,breakable,
  colback=white,colframe=#2,coltitle=white,
  fonttitle=\bfseries,
  attach boxed title to top left={xshift=8pt,yshift=-\tcboxedtitleheight/2},
  boxed title style={colback=#2,arc=2pt,boxrule=0pt},
  boxrule=0.9pt,arc=2pt,
  left=8pt,right=8pt,top=8pt,bottom=6pt,
  title={#3},#1}

\newcounter{methode}
\newenvironment{definitionbox}[1][Définition]%
  {\begin{boitecours}[phantom={\signet{2}{#1}}]{coulDef}{#1}}{\end{boitecours}}
\newenvironment{proprietebox}[1][Propriété]%
  {\begin{boitecours}[phantom={\signet{2}{#1}}]{coulProp}{#1}}{\end{boitecours}}
\newenvironment{methodebox}[1][Méthode]{\stepcounter{methode}%
  \begin{boitecours}[phantom={\signet[methode-\arabic{methode}]{2}{#1}}]{coulMeth}{#1}}%
  {\end{boitecours}}

\newcommand{\etape}[1]{\par\smallskip\textbf{\color{coulMeth}#1}\par\nobreak\smallskip}
\newcommand{\enonce}[1]{\emph{#1}\par\smallskip}
\newcommand{\reponse}[1]{\par\smallskip
  {\footnotesize\itshape\color{black!55}Réponses : #1}}
\newcommand{\demo}[1]{\textbf{\color{coulMeth}Démonstration.} #1}
\newcommand{\afaire}[1]{%
  \par\smallskip
  \noindent{\small\color{coulExo}$\blacktriangleright$\ \textbf{À faire :}
    fiche d'exercices du chapitre, #1.}\par\smallskip}

\newenvironment{calculs}{\setlength{\jot}{5pt}\@fleqntrue\@mathmargin\retraitcalculs
  \setlength{\abovedisplayskip}{4pt}\setlength{\belowdisplayskip}{4pt}%
  \setlength{\abovedisplayshortskip}{4pt}\setlength{\belowdisplayshortskip}{4pt}%
  \csname alignat*\endcsname{2}}%
  {\csname endalignat*\endcsname}
\newcommand{\rem}[1]{\text{\small\itshape\color{black!60}#1}}
\newcommand{\explique}[2]{\par\smallskip\noindent
  \begin{minipage}[t]{0.57\linewidth}\raggedright #1\end{minipage}\hfill
  \begin{minipage}[t]{0.40\linewidth}\raggedright\leavevmode
    {\small\itshape\color{black!60}#2}\end{minipage}\par}
\newcommand{\cas}[1]{\par\smallskip\noindent\textbf{\color{coulExo}#1}\nobreak}
\newcommand{\calc}[1]{\par\smallskip\textbf{\color{coulExo}#1}\par\nobreak}

\newtcolorbox{remarquebox}[1][Remarque]{%
  enhanced,breakable,colback=white,colframe=coulRem,
  boxrule=0pt,leftrule=2.5pt,arc=0pt,outer arc=0pt,
  left=8pt,right=6pt,top=5pt,bottom=5pt,
  before upper={\textbf{\color{coulRem}#1}\par\smallskip}}

\newtcolorbox{exemplebox}[1][Exemple]{%
  enhanced,breakable,colback=white,colframe=coulEx,
  boxrule=0pt,leftrule=2.5pt,arc=0pt,outer arc=0pt,
  left=8pt,right=6pt,top=5pt,bottom=5pt,
  before upper={\textbf{\color{coulEx}#1}\par\smallskip}}

\pgfplotsset{
  repere/.style={
    axis lines=middle,
    axis line style={-{Stealth[length=2.4mm]},thick},
    every axis x label/.style={at={(ticklabel* cs:1.0)},anchor=west},
    every axis y label/.style={at={(ticklabel* cs:1.0)},anchor=south},
    label style={font=\footnotesize},
    tick label style={font=\scriptsize},
    grid=both,
    major grid style={grille},
    minor tick num=0,
    tick align=inside,
    clip mode=individual,
  },
}
\tikzset{
  courbe/.style={coulCourbe,line width=1.1pt,smooth},
  courbeB/.style={coulCourbeB,line width=1.1pt,smooth},
  courbeC/.style={coulCourbeC,line width=1.1pt,smooth},
}

\newlist{etapes}{enumerate}{1}
\setlist[etapes]{label=\textbf{\arabic*.},leftmargin=*,itemsep=1pt,topsep=2pt}

\newlist{expressions}{enumerate}{1}
\setlist[expressions]{label={\Alph*\ $=$},leftmargin=*,itemsep=3pt,topsep=3pt}

\newcounter{exo}
\renewcommand{\theexo}{\ifnum\value{exo}<10 0\fi\arabic{exo}}
\newtcolorbox{exercice}[1][]{%
  enhanced,breakable,
  colback=white,colframe=coulExo,
  boxrule=0.9pt,arc=2pt,
  left=8pt,right=8pt,top=8pt,bottom=6pt,
  coltitle=white,fonttitle=\bfseries,
  attach boxed title to top left={xshift=8pt,yshift=-\tcboxedtitleheight/2},
  boxed title style={colback=coulExo,arc=2pt,boxrule=0pt},
  phantom={\signet{2}{Exercice \arabic{exo}}},
  title={Exercice \theexo\ifx\relax#1\relax\else{} --- #1\fi}}
\newenvironment{exo}[1][\relax]{\refstepcounter{exo}\begin{exercice}[#1]}{\end{exercice}}

  \RequirePackage{comment}
  \excludecomment{correction}
  \renewcommand{\reponse}[1]{}
  \newcommand{\autableau}{\hfill{\footnotesize\itshape\color{black!45}(au tableau)}}
  \newcommand{\etapeexemples}{%
    \par\smallskip\textbf{\color{coulMeth}Exemples}\autableau\par\nobreak\smallskip}
  \newcommand{\etapetableau}[1]{%
    \par\smallskip\textbf{\color{coulMeth}#1}\autableau\par\nobreak\smallskip}
  \newtcolorbox{exempletableau}[1][Exemple]{%
    enhanced,breakable,colback=white,colframe=coulEx,
    boxrule=0pt,leftrule=2.5pt,arc=0pt,outer arc=0pt,
    left=8pt,right=6pt,top=4pt,bottom=4pt,
    before skip=5pt,after skip=5pt,
    before upper={\textbf{\color{coulEx}#1}\autableau\par\smallskip}}

\RequirePackage{comment}
\excludecomment{correctionavj}

\newcommand{\coursEntete}{}
\newcommand{\coursNiveau}{}
\newcommand{\coursTitre}{}
\newcommand{\coursSurTitre}{}
\newcommand{\coursSousTitre}{}

\newcommand{\enteteCours}[1]{\renewcommand{\coursEntete}{#1}}
\newcommand{\chapitrecours}[2]{\enteteCours{Chapitre #1 --- #2}}
\newcommand{\niveaucours}[1]{\renewcommand{\coursNiveau}{#1}}
\newcommand{\titrecours}[3][]{%
  \renewcommand{\coursTitre}{#2}%
  \renewcommand{\coursSurTitre}{#1}%
  \renewcommand{\coursSousTitre}{#3}}

\pagestyle{fancy}
\fancyhf{}
\fancyhead[L]{\footnotesize\color{coulPart}\textbf{\coursEntete}}
\fancyhead[R]{\footnotesize\color{coulPart}\textbf{\coursNiveau}}
\fancyfoot[C]{\footnotesize\thepage}
\renewcommand{\headrulewidth}{0.4pt}
\renewcommand{\headrule}{\hbox to\headwidth{\color{coulPart!40}\leaders\hrule height \headrulewidth\hfill}}

\setlength{\parskip}{2pt}

\AtBeginDocument{%
  \begin{center}
    \begin{tcolorbox}[enhanced,width=0.72\linewidth,colback=white,colframe=coulPart,
        boxrule=1.6pt,arc=3pt,halign=center,top=10pt,bottom=10pt]
      \ifx\coursSurTitre\empty
        {\Huge\bfseries\color{coulPart} \coursTitre}\\[4pt]
      \else
        {\Huge\bfseries\color{coulPart} \coursTitre}\\[3pt]
        {\Large\bfseries\color{coulPart} \coursSurTitre}\\[5pt]
      \fi
      {\normalsize \coursSousTitre}%
      \\[2pt]{\small\itshape Version élève}
    \end{tcolorbox}
  \end{center}%
}
\makeatother

\chapitrecours{1}{Calcul littéral}
\niveaucours{1\textsuperscript{re} STI2D}
\titrecours{CALCUL LITTÉRAL}{Cours --- Première STI2D}

\begin{document}

\begin{remarquebox}[Comment utiliser ce cours]
Chaque partie commence par ce qu'il faut \textbf{savoir} (définitions, théorèmes), puis
vient ce qu'il faut \textbf{savoir faire} : une méthode, avec des
\textbf{exemples} faits en classe, puis des exercices
\og\textbf{À vous de jouer}\fg{} à chercher seul.

\smallskip
Sur cette feuille, seuls les \textbf{énoncés} des exemples figurent : la résolution,
signalée par la mention \emph{(au tableau)}, est à rédiger dans le cahier.

\end{remarquebox}

\partiecours{Réduire une expression littérale}

\begin{definitionbox}[Définition 1 --- Réduire]
Des termes \textbf{semblables} ont la même partie littérale : $4x^{2}$ et $-x^{2}$ sont
semblables, $4x^{2}$ et $4x$ ne le sont pas.

\smallskip
\textbf{Réduire une somme}, c'est regrouper les termes semblables en additionnant leurs
coefficients. \textbf{Réduire un produit}, c'est multiplier les nombres entre eux et les
lettres entre elles.
\end{definitionbox}

\begin{remarquebox}
Un produit se réduit toujours : $5\times 2x=10x$. Une somme pas toujours : dans $5+2x$,
un terme est un nombre, l'autre contient $x$ ; ils ne sont pas semblables, et
$5+2x$ reste ainsi.
\end{remarquebox}

\begin{methodebox}[Méthode 1 --- Réduire une somme ou un produit]
\etapeexemples
Réduire les expressions suivantes.
\begin{questions}[label=\alph*)]
  \item $A=4x-7y+3x^{2}-9x+2y$
  \item $B=3x\times 5x^{2}\times 2y$
\end{questions}

\etape{À vous de jouer}
Réduire, si possible :
\[
  C=8n-3n^{2}+5-2n+6n^{2}-9 \qquad D=-4x\times 3x\times 5 \qquad
  E=-5a\times\left(-2a^{2}\right) \qquad F=9+4t
\]

\end{methodebox}

\afaire{exercices 1 à 3}

\partiecours{Développer}

\begin{definitionbox}[Définition 2 --- Développer, factoriser]
\begin{minipage}[t]{0.46\linewidth}
\vspace{2pt}
\textbf{Développer}, c'est transformer un produit en une somme.

\smallskip
\textbf{Factoriser}, c'est transformer une somme en un produit.
\end{minipage}\hfill
\begin{minipage}[t]{0.50\linewidth}
\vspace{0pt}
\centering
\begin{tikzpicture}[>={Stealth[length=3mm]}]
  \node (g) at (0,0) {$2x(x+5)$};
  \node (d) at (3.4,0) {$2x^{2}+10x$};
  \draw[->,line width=1pt,coulEx] (g.north east) to[bend left=32]
        node[above,font=\scriptsize\bfseries,coulEx]{DÉVELOPPER} (d.north west);
  \draw[->,line width=1pt,coulExo] (d.south west) to[bend left=32]
        node[below,font=\scriptsize\bfseries,coulExo]{FACTORISER} (g.south east);
\end{tikzpicture}
\end{minipage}
\end{definitionbox}

\begin{proprietebox}[Propriété 1 --- Distributivité]
Pour tous nombres $a$, $b$, $c$, $d$ et $k$ :
\[
  k(a+b)=ka+kb \qquad\qquad (a+b)(c+d)=ac+ad+bc+bd
\]
Un signe $-$ devant une parenthèse change le signe de chaque terme de la parenthèse :
$-(3-x)=-3+x$.
\end{proprietebox}

\begin{methodebox}[Méthode 2 --- Développer et réduire]
\etapeexemples
Développer et réduire les expressions suivantes.
\begin{questions}[label=\alph*)]
  \item $A=-4(3x-5)$
  \item $B=(3x+2)(x-4)$
  \item $C=3x(2-x)-(x+1)(2x-5)$
\end{questions}

\etape{À vous de jouer}
Développer et réduire :
\[
  D=6(2x-3) \qquad E=-3(5-2x) \qquad F=7x-2(3x+5)
\]
\[
  G=(2x+5)(x+3) \qquad H=(3x-2)(5-x) \qquad I=(x+3)(2x-1)-x(x+4)
\]

\end{methodebox}

\afaire{exercices 4 à 6}

\partiecours{Les trois identités remarquables}

\begin{proprietebox}[Propriété 2 --- Identités remarquables]
Pour tous nombres $a$ et $b$ :
\begin{itemize}
  \item 1\textsuperscript{re} identité remarquable : \quad $(a+b)^{2}=a^{2}+2ab+b^{2}$ ;
  \item 2\textsuperscript{e} identité remarquable : \quad $(a-b)^{2}=a^{2}-2ab+b^{2}$ ;
  \item 3\textsuperscript{e} identité remarquable : \quad $(a+b)(a-b)=a^{2}-b^{2}$.
\end{itemize}
\end{proprietebox}

\begin{remarquebox}
Ce sont des raccourcis de la double distributivité. Attention :
$(a+b)^{2}\ne a^{2}+b^{2}$ ; le \textbf{double produit} $2ab$ ne s'oublie jamais.
\end{remarquebox}

\begin{methodebox}[Méthode 3 --- Développer avec une identité remarquable]
\etapeexemples
Développer et réduire les expressions suivantes.
\begin{questions}[label=\alph*)]
  \item $A=(x+6)^{2}$
  \item $B=(2x-5)^{2}$
  \item $C=(3x+4)(3x-4)$
  \item $D=(x-2)^{2}-(x+3)(3-x)$
\end{questions}

\etape{À vous de jouer}
Développer et réduire :
\[
  E=(x+7)^{2} \qquad F=(3x-2)^{2} \qquad G=(5x+1)(5x-1) \qquad
  H=(x+4)(x-4)-(2x-3)^{2}
\]

\end{methodebox}

\afaire{exercices 7 à 10}

\partiecours{Factoriser}

Pour \textbf{factoriser}, c'est-à-dire écrire une expression comme un produit, on
cherche dans l'ordre :
\begin{itemize}
  \item un \textbf{facteur commun} à tous les termes, nombre, lettre ou parenthèse
        entière : $ka+kb=k(a+b)$ ;
  \item sinon, une \textbf{identité remarquable} lue de droite à gauche :
        $a^{2}+2ab+b^{2}=(a+b)^{2}$, $a^{2}-2ab+b^{2}=(a-b)^{2}$,
        $a^{2}-b^{2}=(a-b)(a+b)$.
\end{itemize}

\begin{methodebox}[Méthode 4 --- Factoriser]
\etapeexemples
Factoriser les expressions suivantes.
\begin{questions}[label=\alph*)]
  \item $A=15x^{2}-20x$
  \item $B=(x+4)(2x-3)+(x+4)(x+5)$
  \item $C=(2x+3)^{2}-(2x+3)(x-1)$
  \item $D=4x^{2}+28x+49$
  \item $E=25x^{2}-36$
  \item $F=16-(x+3)^{2}$
\end{questions}

\etape{À vous de jouer}
Factoriser :
\[
  G=14x^{2}-21x \qquad H=(x-2)(4x+1)-(x-2)(x-6) \qquad I=x^{2}-10x+25
\]
\[
  J=9x^{2}+6x+1 \qquad K=49x^{2}-4 \qquad L=(x-3)^{2}-36
\]

\end{methodebox}

\afaire{exercices 11 à 14}

\partiecours{Résoudre des équations}

\begin{definitionbox}[Définition 3 --- Solutions d'une équation]
\textbf{Résoudre une équation} d'inconnue $x$, c'est trouver toutes les valeurs de $x$
qui rendent l'égalité vraie : ce sont les \textbf{solutions}. On note $S$ l'ensemble des
solutions.
\end{definitionbox}

\begin{proprietebox}[Propriété 3 --- Produit nul]
Un produit est nul si et seulement si l'un au moins de ses facteurs est nul :
si $A\times B=0$, alors $A=0$~ou~$B=0$.
\end{proprietebox}

\begin{proprietebox}[Propriété 4 --- Équation $x^{2}=k$]
\begin{itemize}
  \item Si $k>0$, l'équation $x^{2}=k$ a deux solutions : $-\sqrt{k}$ et $\sqrt{k}$.
  \item Si $k=0$, elle a une seule solution : $0$.
  \item Si $k<0$, elle n'a pas de solution : un carré n'est jamais négatif.
\end{itemize}
\end{proprietebox}

\begin{methodebox}[Méthode 5 --- Résoudre une équation]
\etapeexemples
Résoudre dans $\R$ les équations suivantes.
\begin{questions}[label=\alph*)]
  \item $6x-5=2x+7$
  \item $2(x-3)=-(x+4)+5$
  \item $(2x+8)(5-3x)=0$
  \item $(2x-1)(x+4)-(2x-1)(3x-2)=0$
  \item $2x^{2}-7x=0$
  \item $x^{2}=25$ ; $x^{2}=-3$
\end{questions}

\etape{À vous de jouer}
Résoudre dans $\R$ :
\[
  5x-8=2x+7 \qquad 3(x+2)=4(x-1)+1 \qquad (3x-6)(x+5)=0
\]
\[
  (x+3)(2x-1)+(x+3)(x-8)=0 \qquad 5x^{2}+2x=0 \qquad x^{2}=81
\]

\end{methodebox}

\afaire{exercices 15 à 18}

\partiecours{Mise en équation}

\begin{methodebox}[Méthode 6 --- Mettre un problème en équation]
\etapeexemples
Un fablab imprime un lot de $150$ pièces avec deux imprimantes 3D. L'imprimante A
consomme $12$~Wh par pièce, l'imprimante B $18$~Wh. Le lot a consommé $2\,220$~Wh en
tout. Combien de pièces chaque imprimante a-t-elle produites ?

\etape{À vous de jouer}
\begin{questions}
  \item Une entreprise a livré $540$ colis en deux jours ; le second jour, elle en a
        livré $60$ de plus que le premier. Combien de colis a-t-elle livrés chaque jour ?
  \item Une plaque rectangulaire a pour longueur $x+5$ et pour largeur $x-2$ (en cm).
        Son périmètre mesure $34$~cm. Trouver $x$, puis les dimensions de la plaque.
\end{questions}

\end{methodebox}

\afaire{exercices 19 à 21}

\partiecours{Résoudre des inéquations}

\begin{definitionbox}[Définition 4 --- Inéquation]
Une \textbf{inéquation} est une inégalité ($<$, $\le$, $>$ ou $\ge$) qui contient une
inconnue $x$. La \textbf{résoudre}, c'est trouver toutes les valeurs de $x$ qui la
vérifient ; ces solutions forment en général un \textbf{intervalle}.
\end{definitionbox}

\begin{proprietebox}[Propriété 5 --- Changer le sens]
On résout une inéquation comme une équation, avec une règle de plus :
\begin{center}
\fbox{\parbox{0.86\linewidth}{\centering
  \textbf{multiplier ou diviser les deux membres par un nombre strictement négatif
  \emph{change le sens} de l'inégalité.}}}
\end{center}
Par exemple, en divisant $-3x<6$ par $-3$ : $\;x>-2$.
\end{proprietebox}

\begin{methodebox}[Méthode 7 --- Résoudre une inéquation du premier degré]
\etapeexemples
Résoudre dans $\R$ les inéquations suivantes, puis représenter les solutions sur une
droite graduée.
\begin{questions}[label=\alph*)]
  \item $3x+2<7-2x$
  \item $3(x-2)\le 5x-1$
\end{questions}

\etape{À vous de jouer}
Résoudre dans $\R$ et donner les solutions sous forme d'intervalle :
\[
  2x-7>5 \qquad -3x+4\ge 10 \qquad 2x-3<6x+5
\]

\end{methodebox}

\afaire{exercices 22 à 24}

\partiecours{Tableaux de signes}

\begin{proprietebox}[Propriété 6 --- Signe de $ax+b$ (avec $a\ne 0$)]
$ax+b$ s'annule pour $x=-\dfrac{b}{a}$ ; \textbf{après} cette valeur, il a le signe de
$a$, \textbf{avant}, le signe contraire.
\begin{center}
\begin{minipage}{0.47\linewidth}\centering
{\small $a>0$ : on commence par $-$}\par\smallskip
\begin{tikzpicture}[scale=0.9]
  \tkzTabInit[lgt=1.6,espcl=2.2]{$x$/1,$ax+b$/1}{$-\infty$,$-\frac{b}{a}$,$+\infty$}
  \tkzTabLine{,-,z,+,}
\end{tikzpicture}
\end{minipage}\hfill
\begin{minipage}{0.47\linewidth}\centering
{\small $a<0$ : on commence par $+$}\par\smallskip
\begin{tikzpicture}[scale=0.9]
  \tkzTabInit[lgt=1.6,espcl=2.2]{$x$/1,$ax+b$/1}{$-\infty$,$-\frac{b}{a}$,$+\infty$}
  \tkzTabLine{,+,z,-,}
\end{tikzpicture}
\end{minipage}
\end{center}
\end{proprietebox}

\begin{exemplebox}[Exemple --- signe de $5x-10$ et de $-2x+8$]
$5x-10$ s'annule en $2$ et $a=5>0$ : négatif avant $2$, positif après.

\smallskip
$-2x+8$ s'annule en $4$ et $a=-2<0$ : positif avant $4$, négatif après.
\end{exemplebox}

\begin{methodebox}[Méthode 8 --- Signe d'un produit $(ax+b)(cx+d)$]
\etapeexemples
Dresser le tableau de signes de $P(x)=(2x-6)(1-4x)$.

\etape{À vous de jouer}
Dresser le tableau de signes de $Q(x)=(3x+9)(2-x)$, puis dire où $Q(x)$ est positif.

\end{methodebox}

\afaire{exercices 25 à 27}

\partiecours{Inéquations produit}

Pour résoudre une inéquation $A(x)\times B(x)>0$ (ou $\ge 0$, $<0$, $\le 0$), on ne
peut pas isoler $x$ : on dresse le \textbf{tableau de signes} du produit, puis on lit
les solutions sur la dernière ligne. Si l'expression n'est pas un produit, on la
\textbf{factorise d'abord}.

\begin{methodebox}[Méthode 9 --- Résoudre une inéquation produit]
\etapeexemples
Résoudre dans $\R$ les inéquations suivantes.
\begin{questions}[label=\alph*)]
  \item $(2x+5)(3-x)\ge 0$
  \item $x^{2}-25<0$
  \item $(x+2)(x-1)(4-x)\ge 0$
\end{questions}

\etape{À vous de jouer}
Résoudre dans $\R$ :
\[
  (x-4)(x+6)<0 \qquad (3x+1)(5-x)\ge 0 \qquad x^{2}-7x\le 0 \qquad (x+3)(x-2)(5-x)>0
\]

\end{methodebox}

\afaire{exercices 28 à 30}

\partiecours{Inéquations quotient}

\begin{definitionbox}[Définition 5 --- Valeur interdite]
On ne divise pas par $0$ : une valeur qui annule le dénominateur d'un quotient est une
\textbf{valeur interdite}. Elle n'est jamais solution ; dans le tableau de signes, on la
marque d'une \textbf{double barre}.
\end{definitionbox}

\begin{remarquebox}
Un quotient suit la même règle des signes qu'un produit. On ne « fait pas passer » le
dénominateur de l'autre côté : son signe est inconnu, on ne saurait pas s'il faut
changer le sens.
\end{remarquebox}

\begin{methodebox}[Méthode 10 --- Résoudre une inéquation quotient]
\etapeexemples
Résoudre dans $\R$ les inéquations suivantes.
\begin{questions}[label=\alph*)]
  \item $\dfrac{3-9x}{2x-4}\le 0$
  \item $\dfrac{3x+2}{x+1}\le 2$
\end{questions}

\etape{À vous de jouer}
Résoudre dans $\R$, après avoir donné la valeur interdite :
\[
  \dfrac{4-x}{x+2}>0 \qquad \dfrac{x+3}{3x-6}\le 0 \qquad \dfrac{x+5}{x-1}\ge 3
\]

\end{methodebox}

\afaire{exercices 31 à 33}

\partiecours{Problèmes avec des inéquations}

\begin{methodebox}[Méthode 11 --- Résoudre un problème par une inéquation]
\etapeexemples
Un atelier de découpe laser fabrique $x$ pièces par jour, avec $0<x\le 300$. Le coût de
production est $C(x)=6x+840$ et la recette $R(x)=18x$, en euros.
\begin{questions}[label=\alph*)]
  \item À partir de combien de pièces par jour l'atelier fait-il un bénéfice ?
  \item Le coût moyen d'une pièce est $\dfrac{C(x)}{x}$. Combien de pièces faut-il
        produire pour que ce coût moyen ne dépasse pas $10$~€ ?
\end{questions}

\etape{À vous de jouer}
\begin{questions}
  \item Un service d'impression 3D propose deux tarifs : A, $0{,}40$~€ par gramme de
        filament ; B, un abonnement de $15$~€ puis $0{,}25$~€ par gramme. À partir de
        combien de grammes (entiers) le tarif B est-il strictement moins cher ?
  \item Le bénéfice, en euros, d'un atelier qui fabrique $x$ objets est
        $B(x)=(x-30)(90-x)$, avec $0\le x\le 120$. Pour quelles valeurs de $x$
        l'atelier fait-il un bénéfice ?
\end{questions}

\end{methodebox}

\afaire{exercices 34 à 36}

\end{document}
